\documentclass{article}
\usepackage{amsthm,amsmath,amssymb,enumitem}
\newtheorem{problem}{Problem}
\newcommand{\Z}{\mathbb{Z}}
\begin{document}
\begin{problem}
Define sets $A$ and $B$ as follows:
\begin{align*}
A&=\{m \in \Z \mid m=6r+12 \text{ for some } r \in \Z\}\\
B&=\{n \in \Z \mid n=3s \text{ for some } s \in \Z\}\\
\end{align*}

\begin{enumerate}[label=(\alph*)]
\item Prove that $A \subseteq B$.
\item Disprove that $B \subseteq A$.
\end{enumerate}
\end{problem}
\textbf{Solution}
\begin{enumerate}[label=(\alph*)]
\item Let $x \in A$. Then $x=6r+12$ for some integer $r$.
Observe that
\begin{align*}
x&=6r+12\\
&=3(2r+4)
\end{align*}

Let $s=2r+4$. Then, $x=3s$. Hence, $x \in B$, and $A \subseteq B$.
\item Let $x \in B$. We shall prove $x \notin A$ for some cases. We have $x=3s$
for some integer $s$. For $x$ to be in $A$, we would need an integer $r$ such that
$x=6r+12$. Then,
\begin{align*}
3s&=6r+12\\
s&=2r+4
\end{align*}

For $r=1$, $s=6$, and for $r=2$, $s=8$. Hence, by choosing $s=7$, we have
$r=\frac32$ which is not an integer. Therefore, $x \notin A$ for particular cases;
hence, $B \nsubseteq A$.
\end{enumerate}
\end{document}
